%% Front matter in the official Elsevier cas-sc syntax (MSSP guide for authors).
\shorttitle{Why RUL models sometimes predict the same value for every engine}
\shortauthors{C. Zhang, Y. Zha}

\title[mode=title]{Why RUL models sometimes predict the same value for every
engine: an audit of 161 training runs and condition-based design rules}

\author[1]{Chi Zhang}[]
\author[1]{Yuguo Zha}[]
\cormark[1]
\cortext[1]{Corresponding author. E-mail: hugozha@nchu.edu.cn}

\affiliation[1]{organization={Department of Biomedical Engineering, College
of Instrument Science and Optoelectronic Engineering, Nanchang Hangkong
University}, city={Nanchang}, postcode={330063}, state={Jiangxi},
country={China}}

\begin{abstract}
Maintenance planning rests on estimating an aero-engine's remaining useful
life (RUL) in cycles. Models combining convolutional, recurrent and
attention layers are standard on NASA's C-MAPSS benchmark, yet each
configuration is usually trained once and judged by a single test error.
We audited 161 runs to ask whether the same configuration would train again
under a different random seed. On the two subsets with six operating
conditions, 37.2\% of FD002 runs and 24.2\% of FD004 runs ended by printing
one fixed number for every engine, the mean of the training labels; none of
81 single-condition runs did. Loss falls normally, so this reads as poor
tuning, although the score used in service is 15--250$\times$ worse.
Checkpoint probes and per-epoch records show why. Every multi-condition run
first spends 16--34 epochs outputting that average, its features recording
operating condition rather than wear and its attention flat; training
improves only if gradients in the convolutional branch jump about a
hundredfold before early stopping. Five checks spot the outcome from a log
or checkpoint, one a formula predicting its test error to 0.01 cycles. We
therefore set window length and depth by condition count, keep several
branches able to carry the degradation signal, and apply three cheap checks
after training. In a five-seed test all 15 runs on FD001/FD002/FD003 trained
reliably (RMSE 13.66, 19.43, 16.09); on FD004 six of eight runs still
failed, the two that worked agreeing at 19.53 $\pm$ 0.07, which the checks
catch in time for a clean retrain. This is not a shortfall in accuracy but
a run that quietly stops learning, and one RMSE cannot show whether a
configuration is dependable. For safety-critical prognostics, the chance
of failing on retraining belongs beside the error; it has rarely been
measured, and this paper gives a way to measure it.
\end{abstract}

\begin{keywords}
Remaining useful life prediction \sep aero-engine maintenance \sep training
reliability \sep constant predictions \sep operating conditions \sep
attention entropy \sep C-MAPSS \sep hybrid deep networks
\end{keywords}

\maketitle
